\section{Related Work}
\label{sec:related}


\noindent\textbf{Tokenization for language, text, and audio.}
Tokenization is a key component of training pipelines for modern transformer-based autoregressive sequence models, and the choice of tokenization approach can have significant impact on model training and downstream performance~\cite{radford2019language}. While there are multiple works exploring the training of ``tokenization-free'' language models \citep{gillick2016bytelm, meta_blt} that directly operate on bit streams, most language models today rely on a text tokenization stage prior to training. 
A common approach is byte pair encoding~\citep{gage1994new, radford2019language}, which compresses input text by merging frequently occurring token sequences into new tokens. For images, \emph{learned} compression schemes present an effective approach: input images can be represented as ``soft tokens'' produced by a pre-trained vision encoder~\citep{liu2023llava}, and full autoregressive image input-output can be achieved with a vector-quantizing autoencoder~\cite{esser2020taming, vqvae}. Similar approaches can be extended to the video domain~\cite{yu2023magvit}. In audio generation and speech synthesis, which share the time-series structure of action prediction, state-of-the-art models typically encode time-series audio data using either frequency-domain spectrogram images~\citep{gong21b_interspeech} or using learned vector quantizers~\citep{zeghidour2021soundstreamendtoendneuralaudio}.

\noindent\textbf{Vision-language-action models.}
Recently, multiple works have developed \emph{generalist} robot policies~\citep{brohan2022rt,octo_2023,bharadhwaj2023roboagent,rt22023arxiv,Doshi24-crossformer,kim2024openvla,wangscaling,cheang2024gr2generativevideolanguageactionmodel}
that are trained on increasingly large robot learning datasets~\citep{open_x_embodiment_rt_x_2023,khazatsky2024droid,walke2023bridgedata,fang2024rh20t,mandlekar2018roboturk,jiang2024dexmimicgen}. 
One promising approach for training generalist policies are vision-language-action models (VLAs; ~\citep{rt22023arxiv,collaboration2023open,kim2024openvla,Zawalski24-ecot,black2024pi_0,wen2024tinyvlafastdataefficientvisionlanguageaction,zheng2024tracevla,zhen20243d,cheng2024navila,cheang2024gr2generativevideolanguageactionmodel}).
VLAs fine-tune vision-language models, that are pre-trained on internet-scale image and text data, for robot control. This has multiple benefits: using large vision-language model backbones, with billions of parameters, provides policies with the necessary expressivity for fitting large robot datasets. Reusing weights pre-trained on internet-scale datasets also improves the ability of VLAs to follow diverse language commands and generalize, e.g., to new objects and scene backgrounds~\citep{rt22023arxiv,kim2024openvla,Zawalski24-ecot,wen2024tinyvlafastdataefficientvisionlanguageaction,jones2025beyond}. 
Most VLA models today are confined to rather simple, low-frequency control tasks, particularly models that use the most common autoregressive VLA design~\citep{rt22023arxiv, kim2024openvla}. We show that this is a direct consequence of the \emph{action tokenization} schemes employed by these models, which make training on dexterous tasks challenging. We introduce a new action tokenization approach that allows us to train the first autoregressive VLAs on dexterous and high-frequency robot data.

\noindent\textbf{Action representations for VLA training.}
Prior works have explored various action parameterizations for training robot policies, including VLAs.
One line of work uses ``semantic'' action representations like language sub-tasks~\citep{driess2023palm,saycan2022arxiv,belkhale2024rthactionhierarchiesusing}, 
or keypoints~\citep{nasirianypivot,huang2024rekep,fangandliu2024moka,dipalo2024kat}. 
Such approaches can often learn from few examples or even perform tasks \emph{zero-shot} without any robot examples~\citep{nasirianypivot,huang2024rekep,fangandliu2024moka}, but require hand-designed low-level controllers for task execution, limiting their generality. An alternative approach directly trains VLAs to output low-level robot control commands given image and language instruction inputs. The most common design directly embeds actions into discrete tokens, that can be generated with standard autoregressive sequence models, like any popular vision-language model. Existing approaches map from continuous robot actions to discrete action tokens using a simple per-dimension, per-timestep binning scheme~\citep{brohan2022rt,rt22023arxiv,kim2024openvla}.
We find that this scheme struggles to scale to high-frequency robot control tasks. We propose a new tokenization scheme for robot actions, based on time-series compression techniques, that allows us to train autoregressive VLAs on high-frequency data. A number of works have also proposed alternatives to tokenization, for example by using regression heads or introducing new weights for diffusion decoding~\citep{Doshi24-crossformer,black2024pi_0,lee2024behavior,wen2024tinyvlafastdataefficientvisionlanguageaction}. In comparison, our approach does not require modifications of the underlying pre-trained transformer model, can easily be applied to any pre-trained autoregressive transformer model, and achieves competitive performance to state-of-the-art diffusion-based VLAs~\citep{black2024pi_0} across many tasks, while being significantly more compute efficient to train. 

Another set of related work explores \textit{vector-quantized} action representations~\citep{lee2024behavior, belkhale2024minivla, mete2024questselfsupervisedskillabstractions}. Such approaches train a vector-quantized encoder-decoder network, for which reconstruction quality can be sensitive to hyperparameter choices and structure~\cite{yu2023magvit}. We find that these methods perform well at coarse, low-fidelity reconstruction tasks, but fail on high-frequency tasks when fine-grained control is required. In comparison, our {\ModelAcronym} tokenization scheme has few hyperparameters and can reconstruct actions with high precision while offering strong compression properties.



